\documentclass{situationist-preprint}

\title{Gavagai Means Rabbit}
\msid{LNG-06-a}
\msfamily{LNG-06}
\msdiscipline{Philosophy of Language}
\mstype{Principal manuscript}
\mskeywords{radical translation; indeterminacy; Quine; gavagai; stimulus meaning; charity; eligibility; Evans}

\begin{document}
\maketitle

\begin{abstract}
Quine argued that a linguist translating a wholly unknown language could equally well translate the native's \enquote{gavagai} as \enquote{rabbit}, \enquote{undetached rabbit part} or \enquote{rabbit stage}, since all three have the same stimulus meaning, and that the choice is not merely underdetermined by evidence but has no fact to answer to. We argue that radical translation is constrained by more than stimulus meaning, and that the constraints select \enquote{rabbit}. A manual of translation must treat the native's constructions uniformly and compositionally, assign eligible rather than gerrymandered values, and attribute to the native beliefs and interests that are intelligible given her circumstances. Following Evans, we show that once the manual must also translate the native's predication, identity and counting, the rival translations of \enquote{gavagai} survive only by deviant translations of these constructions, which violate the constraints of uniformity and eligibility. Quine's global proxy functions fail the eligibility constraint in the same way. The rabbit manual is the unique manual that satisfies all the constraints.
\end{abstract}

\section{The problem}

In the second chapter of \emph{Word and Object}, \citet{Quine1960} imagines a linguist confronted with a language unrelated to any she knows. A rabbit scurries by, the native says \enquote{gavagai}, and the linguist tentatively writes \enquote{rabbit}. She tests the hypothesis by querying \enquote{gavagai?} in various circumstances and recording assent and dissent. The \emph{stimulus meaning} of a sentence for a speaker is the pair of classes of stimulations that would prompt her assent and her dissent. The stimulus meaning of the one-word sentence \enquote{gavagai} may be the same as that of \enquote{rabbit}. But whenever a rabbit is present, so is an undetached rabbit part and a temporal stage of a rabbit, so the stimulus meanings of \enquote{rabbit}, \enquote{undetached rabbit part} and \enquote{rabbit stage} coincide. No amount of pointing and querying discriminates them. Nor can the linguist ask \enquote{is this gavagai the same as that one?}, since the native's identity predicate can be translated in compensating ways \citep{Quine1960,Quine1969}. He concludes that translation is indeterminate: incompatible manuals fit all possible evidence, and there is no fact of the matter which is right.

Three lines of reply are prominent. \citet{Evans1975} argued that the compensating translations break down once the manual must account for the native's predication: the translation of \enquote{gavagai} cannot be adjusted independently of the translations of adjectives and of the construction that combines them. \citet{Grandy1973} proposed a principle of humanity, that the translator should attribute to the native a pattern of beliefs and desires as similar to our own as the evidence allows, and \citet{Lewis1974} made the attribution of attitudes part of the interpreter's task. \citet{Lewis1983,Lewis1984} argued that interpretation is constrained by the eligibility of the values assigned, so that natural properties are preferred to gerrymandered ones; \citet{Williams2007} examined how far eligibility settles questions of reference. Children learning words assume by default that a novel noun applied to an object names the whole object \citep{Markman1990}, which suggests that the native, if she learned \enquote{gavagai} as children learn nouns, means rabbit; our argument does not depend on this.

We argue that the constraints jointly select \enquote{rabbit}.

\section{Constraints on translation}

\begin{definition}[Manual]\label{d:manual}
A \emph{manual} of translation is a finite recursive specification that maps the native's words and constructions into the translator's language, and so maps every native sentence to a sentence of the translator's.
\end{definition}

A manual is correct if it is the manual that best satisfies the following constraints; if two manuals satisfied them equally, the choice between them would be indeterminate.

\begin{principle}[Stimulus fit]\label{p:stim}
The manual maps occasion sentences to sentences with the same stimulus meaning.
\end{principle}

\begin{principle}[Uniformity]\label{p:unif}
The manual translates each native word and each native construction by a single item, and the translation of a complex expression is fixed by the translations of its parts and of its construction.
\end{principle}

\begin{principle}[Eligibility]\label{p:elig}
Other things equal, the manual assigns natural values (objects, properties and relations that are not gerrymandered) to the native's expressions.
\end{principle}

\begin{principle}[Humanity]\label{p:hum}
Other things equal, the manual attributes to the native beliefs and interests intelligible given her circumstances.
\end{principle}

Principle~\ref{p:stim} is Quine's. Principle~\ref{p:unif} is the interpreter's version of the compositional truth-conditional constraint of LNG-01. Principle~\ref{p:elig} is Lewis's, and is the same constraint that, on EPI-04, makes \enquote{green} rather than \enquote{grue} projectible. Principle~\ref{p:hum} is Grandy's. Only the first is behavioral; the others concern the structure of the manual and the attitudes it attributes.

\section{The argument}

\begin{lemma}[Predication]\label{l:pred}
A manual that translates \enquote{gavagai} as \enquote{undetached rabbit part} or \enquote{rabbit stage}, and satisfies Principle~\ref{p:stim} for the native's compound sentences, must translate her adjectival construction or her identity predicate non-uniformly or by a gerrymandered relation.
\end{lemma}

\begin{proof}
Suppose the native has a color word that the manual translates as \enquote{white}, a construction that combines it with \enquote{gavagai}, and an identity construction, and that she assents to the compound \enquote{white gavagai} when shown a rabbit that is white all over and dissents when shown one with a white tail and a brown body. On the rabbit manual, with predication translated as predication of the same object, the compound means \enquote{white rabbit}, and her responses fit. On the part manual with the same translation of predication, the compound means \enquote{white undetached rabbit part}, which is true of the brown rabbit's tail; her dissent does not fit. To restore fit, the part manual must translate the construction not as predication but as \enquote{is a part of a white rabbit}, a relation holding between a part and the whole to which it belongs, and must do the same for every adjective; and it must translate the identity construction by \enquote{is part of the same rabbit as} \citep{Evans1975}. The stage manual must similarly translate predication as \enquote{is a stage of a rabbit that is white throughout} and identity as \enquote{is a stage of the same rabbit as}.
\end{proof}

\begin{lemma}[Cost]\label{l:cost}
The compensating translations of Lemma~\ref{l:pred} violate Principle~\ref{p:unif} or Principle~\ref{p:elig}.
\end{lemma}

\begin{proof}
The construction that combines adjectives with nouns also combines them with words the manual translates as names of whole objects, such as the native's words for persons, rivers and huts. Either the manual translates the construction as \enquote{part of a white \ldots} only when the noun is \enquote{gavagai} and its kin, and as plain predication elsewhere, which violates Principle~\ref{p:unif}; or it translates it as \enquote{part of a white \ldots} throughout, and so must translate all the native's nouns as part-terms, every identity statement as one of being part of the same whole, and every counting construction as counting wholes by way of their parts. The resulting manual assigns to the native's most basic constructions relations defined by detour through the very wholes the rabbit manual names directly. These relations are less natural than predication and identity, and the manual violates Principle~\ref{p:elig}. The same holds for stages, with persistence through time in place of parthood.
\end{proof}

\begin{theorem}\label{t:main}
Of the manuals that satisfy Principle~\ref{p:stim}, only the rabbit manual satisfies Principles~\ref{p:unif}, \ref{p:elig} and~\ref{p:hum}. So \enquote{gavagai} means rabbit.
\end{theorem}

\begin{proof}
By Lemmas~\ref{l:pred} and~\ref{l:cost}, the part and stage manuals violate Principle~\ref{p:unif} or~\ref{p:elig}. Translating \enquote{gavagai} as \enquote{rabbithood} (as a mass or abstract term) fails in the same way, since the native counts gavagai, and counting must then be rendered as counting the manifestations of a universal. The rabbit manual translates predication as predication and identity as identity, assigns natural values throughout, and attributes to the native an interest in rabbits, animals she hunts, eats and avoids, rather than in their undetached parts; it satisfies Principle~\ref{p:hum} where the rivals attribute an interest that nothing in her circumstances explains.
\end{proof}

\section{Objections and replies}

\begin{objection}[Proxy functions]
Quine's later argument does not depend on rival translations of particular words. Given any manual, apply a proxy function, a one-to-one mapping of objects to objects, to the values of every term, and reinterpret every predicate accordingly. Every sentence keeps its truth value, so the reinterpreted manual fits all evidence \citep{Quine1969}.
\end{objection}

\begin{reply}
A global proxy reinterpretation preserves truth values, and so satisfies Principles~\ref{p:stim} and~\ref{p:unif}, but it assigns to the native's predicates the images of natural properties under an arbitrary permutation, which are gerrymandered. It violates Principle~\ref{p:elig}, and does so maximally, since it does so for every predicate. This is the reply \citet{Lewis1984} gave to the parallel argument from model theory. Eligibility may not settle every question of reference \citep{Williams2007}, but it settles this one.
\end{reply}

\begin{objection}[Begging the question]
Eligibility and humanity are not evidence. They are the translator's preferences, and using them to settle what the native means assumes there is a fact to settle.
\end{objection}

\begin{reply}
The constraints are not evidence in the sense of observations, but they are not preferences. They are constraints on what it is for a manual to be correct, and they are the constraints on which every interpretation, including the interpretation of our own language, proceeds. Quine's argument from indeterminacy assumes that only stimulus meaning constrains translation; if the further constraints are legitimate, the argument does not start. That they are legitimate is supported by the fact that the native means something by \enquote{gavagai}, a primitive fact on LNG-05, and the constraints are what make that fact knowable by others, as LNG-07 requires of any rule a language can have.
\end{reply}

\begin{objection}[Underdetermination]
At most the argument shows that the rabbit manual is better supported. Underdetermination of theory by evidence is the normal condition of science, and Quine claimed only that translation is underdetermined in a further way.
\end{objection}

\begin{reply}
Quine claimed that translation, unlike physics, is indeterminate: that there is no fact for the best manual to get right. Our conclusion is that the constraints select one manual, so that translation is at most underdetermined by stimulus meaning in the way that physics is underdetermined by observation. On SCI-04, physics is underdetermined only by empirically equivalent rivals that lack independent motivation; the part and stage manuals are rivals of exactly this kind.
\end{reply}

\section{What the result does not establish}

The argument does not establish that every question of translation has a determinate answer, only that Quine's paradigm case does. Nor does it reduce meaning to the constraints. On MND-12, no condition stated in non-intentional terms fixes content, and Principle~\ref{p:hum} is not stated in non-intentional terms. LNG-16 extends the result from translation to reference.

\section{Conclusion}

Stimulus meaning does not distinguish \enquote{rabbit} from \enquote{undetached rabbit part} and \enquote{rabbit stage}. But a manual must translate the native's predication, identity and counting as well as her nouns, and the rivals to \enquote{rabbit} survive only by translating these constructions non-uniformly or by gerrymandered relations, while attributing to the native interests nothing explains. Under the constraints of uniformity, eligibility and humanity, \enquote{gavagai} means rabbit.

\begin{thebibliography}{99}\small
\bibitem[Evans(1975)]{Evans1975} Evans, G. (1975). Identity and predication. \emph{Journal of Philosophy}, 72(13), 343--363.
\bibitem[Grandy(1973)]{Grandy1973} Grandy, R. (1973). Reference, meaning, and belief. \emph{Journal of Philosophy}, 70(14), 439--452.
\bibitem[Lewis(1974)]{Lewis1974} Lewis, D. (1974). Radical interpretation. \emph{Synthese}, 27, 331--344.
\bibitem[Lewis(1983)]{Lewis1983} Lewis, D. (1983). New work for a theory of universals. \emph{Australasian Journal of Philosophy}, 61(4), 343--377.
\bibitem[Lewis(1984)]{Lewis1984} Lewis, D. (1984). Putnam's paradox. \emph{Australasian Journal of Philosophy}, 62(3), 221--236.
\bibitem[Markman(1990)]{Markman1990} Markman, E.~M. (1990). Constraints children place on word meanings. \emph{Cognitive Science}, 14(1), 57--77.
\bibitem[Quine(1960)]{Quine1960} Quine, W.~V.~O. (1960). \emph{Word and Object}. MIT Press.
\bibitem[Quine(1969)]{Quine1969} Quine, W.~V.~O. (1969). Ontological relativity. In \emph{Ontological Relativity and Other Essays}. Columbia University Press.
\bibitem[Williams(2007)]{Williams2007} Williams, J.~R.~G. (2007). Eligibility and inscrutability. \emph{Philosophical Review}, 116(3), 361--399.
\end{thebibliography}

\end{document}
